\SourceLeaf{289}%
\begin{center}
\begin{tabular}{lr}
Chasseral-Suchet&\(+0''04536\)\\
Chasseral-Berra&\(-0{,}00966\)\\
Suchet-Chasseral&\(+0{,}06221\)\\
Suchet-Berra&\(+0{,}01014\)\\
Berra-Suchet&\(-0{,}04717\)\\
Berra-Chasseral&\(-0{,}06039\)
\end{tabular}
\end{center}
also auch hier ohne Einfluss auf die Rechnung, die in dem angeführten Werk auf Zehntelsekunden geführt wird.
\Article{15}
Die angegebenen Vorschriften lösen die in den Artikeln 12 und 13 behandelte Aufgabe mit einer Genauigkeit, die für die Anwendung mehr als ausreicht. Dennoch lohnt es sich, die Genauigkeit in einem besonderen Fall um eine weitere Ordnung zu steigern. Für den gleichmäßigen Abschluss einer Untersuchung, die später mitgeteilt werden soll, ist das sogar notwendig. In diesem besonderen Fall muss die Linie \(N\) den Normalparallelkreis in einem zwischen \(F\) und \(G\) gelegenen Punkt \(H\) treffen. Dann ist es günstiger, den Ursprung von \(x\) von \(F\), wo er oben lag, nach \(H\) zu verlegen. Entwickelt man \(l\) und \(m\) in Reihen nach aufsteigenden Potenzen von \(x\), so fallen dadurch in der ersten Reihe das erste und zweite Glied und in der zweiten Reihe das zweite und dritte Glied weg. Die Reihen haben also die folgende Form
\[
\begin{aligned}
l&=\lambda xx+\lambda'x^3+\text{usw.}\\
m&=1+\mu x^3+\mu'x^4+\text{usw.}
\end{aligned}
\]
\indent Für unseren Zweck ist von den Koeffizienten in diesen Reihen nur \(\lambda\) erforderlich. Der folgende Ausdruck dafür lässt sich leicht aus der in Artikel 9 für \(\log m\) angegebenen Formel und den Entwicklungen aus Artikel 13 herleiten
\[
\lambda=-\frac{2ee\cos P\sin P\sin\chi\cos\chi^2}{\cos\varphi\cos\theta}
\]
Dabei behalten \(e,P,\varphi,\theta\) die oben erläuterten Bedeutungen. Für \(\chi\) ist das im Punkt \(H\) geltende Azimut des Bogens \(N\) einzusetzen.

Verwendet man die obigen Reihen bei der Integration der Gleichungen
